\newpage

% =========================================================
% ВАРИАНТ 31
% =========================================================

\section*{Вариант 31}

Четырёхугольник \(ABCD\) вписан в окружность.
Угол \(ABC\) равен \(122^\circ\), угол \(ABD\) равен \(36^\circ\).
Найдите угол \(CAD\). Ответ дайте в градусах.

\begin{center}
\begin{tikzpicture}[
  scale=1.08,
  line cap=round,
  line join=round,
  every node/.style={font=\small}
]
  \def\R{2.45}

  \coordinate (O) at (0,0);

  % Положение точек обеспечивает:
  % угол ABC = 122 градусов;
  % угол ABD = 36 градусов;
  % угол CAD = 86 градусов.
  \coordinate (A) at (200:\R);
  \coordinate (B) at (250:\R);
  \coordinate (C) at (316:\R);
  \coordinate (D) at (128:\R);

  % Окружность
  \draw[line width=0.9pt] (O) circle (\R);

  % Вписанный четырёхугольник
  \draw[line width=0.9pt]
    (A)--(B)--(C)--(D)--cycle;

  % Диагонали, необходимые для указанных углов
  \draw[line width=0.8pt] (A)--(C);
  \draw[line width=0.8pt] (B)--(D);

  % Угол ABD = 36 градусов
  \pic[
    draw,
    line width=0.65pt,
    angle radius=5mm
  ] {angle=A--B--D};

  % Угол ABC = 122 градуса
  \pic[
    draw,
    line width=0.65pt,
    angle radius=9mm
  ] {angle=A--B--C};

  \node at (-0.62,-1.63) {$36^\circ$};
  \node at (0.00,-1.98) {$122^\circ$};

  % Подписи
  \node[below left]  at (A) {$A$};
  \node[below]       at (B) {$B$};
  \node[below right] at (C) {$C$};
  \node[above left]  at (D) {$D$};
\end{tikzpicture}
\end{center}

\[
\boxed{\text{Ответ: }86}
\]


% =========================================================
% ВАРИАНТ 32
% =========================================================

\section*{Вариант 32}

Два угла треугольника равны \(68^\circ\) и \(35^\circ\).
Найдите тупой угол, который образуют высоты треугольника,
выходящие из вершин этих углов.
Ответ дайте в градусах.

\begin{center}
\begin{tikzpicture}[
  scale=1.02,
  line cap=round,
  line join=round,
  every node/.style={font=\small}
]
  % Неравнобедренный остроугольный треугольник:
  % угол A = 68 градусов, угол B = 35 градусов.
  \coordinate (A) at (0,0);
  \coordinate (B) at (7,0);
  \coordinate (C) at (1.544,3.821);

  % Основания высот
  \coordinate (D) at (2.303,3.289);
  \coordinate (E) at (0.982,2.431);

  % Ортоцентр
  \coordinate (H) at (1.544,2.205);

  % Треугольник
  \draw[line width=0.9pt]
    (A)--(B)--(C)--cycle;

  % Высоты AD и BE
  \draw[line width=0.85pt] (A)--(D);
  \draw[line width=0.85pt] (B)--(E);

  % Прямые углы у оснований высот
  \pic[
    draw,
    line width=0.6pt,
    angle radius=3mm
  ] {right angle=A--D--C};

  \pic[
    draw,
    line width=0.6pt,
    angle radius=3mm
  ] {right angle=B--E--C};

  % Данные углы треугольника
  \pic[
    draw,
    line width=0.65pt,
    angle radius=7mm
  ] {angle=B--A--C};

  \pic[
    draw,
    line width=0.65pt,
    angle radius=7mm
  ] {angle=C--B--A};

  \node at (0.78,0.42) {$68^\circ$};
  \node at (5.98,0.34) {$35^\circ$};

  % Искомый тупой угол между высотами
  \pic[
    draw,
    line width=0.7pt,
    angle radius=5.5mm
  ] {angle=A--H--B};

  \node at (1.65,1.50) {$?$};

  \fill (H) circle (1.3pt);

  % Подписи
  \node[below left]  at (A) {$A$};
  \node[below right] at (B) {$B$};
  \node[above]       at (C) {$C$};
  \node[above right] at (D) {$D$};
  \node[left]        at (E) {$E$};
  \node[below]       at (H) {$H$};
\end{tikzpicture}
\end{center}

\[
\boxed{\text{Ответ: }103}
\]


% =========================================================
% ВАРИАНТ 33
% =========================================================

\newpage

\section*{Вариант 33}

Найдите угол \(ACO\), если его сторона \(CA\) касается окружности
с центром \(O\), отрезок \(CO\) пересекает окружность в точке \(B\),
а дуга \(AB\) окружности, заключённая внутри этого угла,
равна \(17^\circ\).
Ответ дайте в градусах.

\begin{center}
\begin{tikzpicture}[
  scale=1.02,
  line cap=round,
  line join=round,
  every node/.style={font=\small}
]
  \def\R{3.2}
  \pgfmathsetmacro{\Cx}{\R/cos(17)}

  \coordinate (O) at (0,0);
  \coordinate (B) at (0:\R);
  \coordinate (A) at (17:\R);
  \coordinate (C) at (\Cx,0);

  % Окружность
  \draw[line width=0.9pt] (O) circle (\R);

  % Радиус, секущая и касательная
  \draw[line width=0.8pt] (O)--(A);
  \draw[line width=0.9pt] (O)--(C);
  \draw[line width=0.9pt] (A)--(C);

  % Выделяем дугу AB
  \draw[line width=1.4pt]
    (0:\R) arc[start angle=0,end angle=17,radius=\R];

  % Радиус перпендикулярен касательной
  \pic[
    draw,
    line width=0.65pt,
    angle radius=3.2mm
  ] {right angle=O--A--C};

  % Искомый угол ACO
  \pic[
    draw,
    line width=0.65pt,
    angle radius=5mm
  ] {angle=A--C--O};

  % Градусная мера дуги
  \node at (8.5:3.72) {$17^\circ$};

  \fill (O) circle (1.2pt);
  \fill (A) circle (1.1pt);
  \fill (B) circle (1.1pt);

  % Подписи
  \node[below left]  at (O) {$O$};
  \node[below]       at (B) {$B$};
  \node[above]       at (A) {$A$};
  \node[right]       at (C) {$C$};
\end{tikzpicture}
\end{center}

\[
\boxed{\text{Ответ: }73}
\]


% =========================================================
% ВАРИАНТ 34
% =========================================================

\section*{Вариант 34}

Угол \(ACO\) равен \(62^\circ\).
Его сторона \(CA\) касается окружности с центром в точке \(O\).
Отрезок \(CO\) пересекает окружность в точке \(B\).
Найдите градусную меру дуги \(AB\) окружности,
заключённой внутри этого угла.
Ответ дайте в градусах.

\begin{center}
\begin{tikzpicture}[
  scale=1.04,
  line cap=round,
  line join=round,
  every node/.style={font=\small}
]
  \def\R{2.9}
  \pgfmathsetmacro{\Cx}{\R/cos(28)}

  \coordinate (O) at (0,0);
  \coordinate (B) at (0:\R);
  \coordinate (A) at (28:\R);
  \coordinate (C) at (\Cx,0);

  % Окружность
  \draw[line width=0.9pt] (O) circle (\R);

  % Радиус, секущая и касательная
  \draw[line width=0.8pt] (O)--(A);
  \draw[line width=0.9pt] (O)--(C);
  \draw[line width=0.9pt] (A)--(C);

  % Искомая дуга AB
  \draw[line width=1.4pt]
    (0:\R) arc[start angle=0,end angle=28,radius=\R];

  % Прямой угол между радиусом и касательной
  \pic[
    draw,
    line width=0.65pt,
    angle radius=3.2mm
  ] {right angle=O--A--C};

  % Угол ACO = 62 градуса
  \pic[
    draw,
    line width=0.65pt,
    angle radius=6mm
  ] {angle=A--C--O};

  \node at (2.69,0.49) {$62^\circ$};

  \fill (O) circle (1.2pt);
  \fill (A) circle (1.1pt);
  \fill (B) circle (1.1pt);

  % Подписи
  \node[below left] at (O) {$O$};
  \node[below]      at (B) {$B$};
  \node[above]      at (A) {$A$};
  \node[right]      at (C) {$C$};
\end{tikzpicture}
\end{center}

\[
\boxed{\text{Ответ: }28}
\]

\end{document}